\documentclass[11pt,a4paper]{article}
\usepackage[margin=27mm]{geometry}
\usepackage{fontspec,unicode-math,amsmath,amsthm,booktabs,longtable,hyperref}
\setmainfont{STIXGeneral}\setmathfont{STIX Two Math}\setmonofont{Menlo}
\newtheorem{theorem}{Theorem}\newtheorem{lemma}[theorem]{Lemma}\newtheorem{proposition}[theorem]{Proposition}
\newcommand{\Aut}{\operatorname{Aut}}\newcommand{\Stab}{\operatorname{Stab}}\newcommand{\cl}{\mathcal C}
\title{Coset recognition and exhaustive maximal-subgroup enumeration}
\author{Common proof for finite-group census applications}
\date{}
\begin{document}\maketitle
\begin{abstract}
Flag-transitive rank-three geometries with distinct maximal object stabilizers can be recognized by three subgroup-generation equations and one exact completion-orbit count. Stabilizer-orbit augmentation then exhausts the admissible unordered triples. We give the full elementary reductions used in the accompanying Mathieu census. The complete finite-group catalogue and exact algorithms remain separate application premises.
\end{abstract}
Throughout this common paper, $G$ is a finite group. This finiteness assumption is required when equality of cardinalities identifies the completion set with a single orbit.

\section{Coset recognition and exhaustive enumeration}
An incidence geometry here is a typed set with symmetric incidence only between different types, with at least one chamber containing every type. Flag transitivity means transitivity on flags of each fixed type set. Residual connectedness means connectedness of every residue of rank at least two, including the whole geometry. We distinguish a specified acting group $G$. Our equivalence permits an automorphism of $G$ and reordering of the types, as in the source's unordered subgroup notation. We do not assert that every abstract isomorphism of the incidence structures normalizes this distinguished group.

\begin{lemma}[Coset realization and repeated stabilizers]\label{lem:realize}
Let $G$ act flag-transitively, and let $(x_i)$ be a base chamber with $H_i=\Stab_G(x_i)$. The map $H_i g\mapsto x_i^g$ identifies the geometry with the right-coset system, whose incidence is nonempty coset intersection. In rank two or three, residual connectedness and maximal proper object stabilizers imply that the $H_i$ are distinct.
\end{lemma}
\begin{proof}
Transitivity on single objects makes each map surjective. The equations $x_i^g=x_i^h$, $gh^{-1}\in H_i$ and $H_i g=H_i h$ are equivalent, proving well-definedness and injectivity. A common coset representative gives a translate of a base incident pair. Conversely, flag transitivity moves the base pair to any incident pair, supplying a common representative. This proves both implications of incidence.

Two equal stabilizers $H$ give a rank-two graph that is a disjoint union of edges indexed by $H\backslash G$. If all three stabilizers coincide, the graph is instead a disjoint union of triangles with the same index set. A proper subgroup has at least two cosets. Finally suppose exactly two rank-three stabilizers equal $H$, while the third is a distinct maximal subgroup $K$. The residue at the base $K$-object is a disjoint union of edges indexed by $(H\cap K)\backslash K$. Distinct maximality implies $K\not\leq H$, so this index is at least two. Each case contradicts the required connectedness.
\end{proof}

For three distinct maximal proper subgroups put
\[
 H_{ij}=H_i\cap H_j,\quad B=H_1\cap H_2\cap H_3,\qquad
 R_{ij}=\{H_i h:h\in H_j\}.
\]
\begin{theorem}[A complete rank-three test]\label{thm:test}
Their coset system is flag-transitive and residually connected if and only if
\begin{equation}\label{eq:test}
 H_i=\langle H_{ij},H_{ik}\rangle\quad(i=1,2,3),\qquad
 |R_{ij}\cap R_{ik}|=[H_{jk}:B]
\end{equation}
for one choice of distinct $i,j,k$ in the cardinality equation. Every accepted system is firm. Two distinct maximal subgroups always give a firm connected flag-transitive rank-two geometry.
\end{theorem}
\begin{proof}
An incident pair of cosets has a common representative, so right multiplication is transitive on each pair type. Fix the base pair $H_j,H_k$. Its possible completions are exactly $R_{ij}\cap R_{ik}$. The pair stabilizer $H_{jk}$ has an orbit of size $[H_{jk}:B]$ containing the base completion $H_i$. Equality says that this orbit is the whole completion set. Consequently all chambers are translates of the base chamber, and every incident pair extends to one. This is equivalent to flag transitivity, since singleton and pair transitivity already hold.

Under flag transitivity the residue at $H_i$ is the bipartite coset graph inside $H_i$ on $H_{ij},H_{ik}$. Its connected component at the base edge consists precisely of the cosets represented by words in those two subgroups: each factor extends a path, and every path gives such a word. Hence that residue is connected precisely when the corresponding generation equation holds. Any two distinct maximal subgroups generate $G$, so the whole coset graph is connected. These are all residues of rank at least two.

If a panel has one completion, then some $H_{jk}=B$. Generation at $H_j$ gives $H_j=\langle H_{ij},B\rangle=H_{ij}$, contradicting distinct maximality. Thus panels contain at least two objects. The rank-two assertions follow from the same path argument and from $[H_i:H_1\cap H_2]>1$ for both $i$.
\end{proof}

\begin{proposition}[Exhaustive augmentation]\label{prop:orbits}
Let $\mathcal M$ be the complete set of maximal subgroups of a finite group, and let $A=\Aut(G)$ act on it. First enumerate $A$-orbits on unordered pairs. For each representative $e$, enumerate $A_e$-orbits on $\mathcal M\setminus e$, attach one representative and apply Theorem~\ref{thm:test}. Fusion under $A$ gives exactly the admissible unordered triple classes.
\end{proposition}
\begin{proof}
Pairs can themselves be enumerated by fixing a representative $H$ of each $A$-orbit and then taking $A_H$-orbits on the remaining members, followed by unordered-pair fusion. Given any triple, choose one of its pairs and move that pair to its representative $e$. An element of the setwise stabilizer $A_e$ then moves its remaining member to a tested representative. Recognition is invariant under these moves. Conversely every accepted triple passes the geometric criterion. Different distinguished pairs can lead to the same triple orbit, so final fusion is necessary and is performed. Invariant size signatures may restrict the transporter searches, but equality of such signatures is never substituted for an exact transporter.

The actual permutation image of $A$ on $\mathcal M$ suffices: every image element lifts to an automorphism and gives exactly the same orbits. An ambient normalizer must first be divided by its conjugation kernel; its order alone is not the automorphism-group order.
\end{proof}


\section*{Attribution and scope}
Coset realization, connectedness by subgroup generation and stabilizer-orbit augmentation are classical. The three complete proofs above are reproduced from the preceding corrected Mathieu proof edition, without using its group-specific census, geometric coordinates or software outputs. Their source application is the classification of Leemans and Rowley cited below. This paper does not establish catalogue completeness or a named-group identification. A complete application must supply those inputs and its full required computation.
\begin{thebibliography}{9}
\bibitem{LR} D. Leemans and P. Rowley, \emph{Rank three residually connected geometries for $M_{22}$, revisited}, Electronic Journal of Combinatorics 17 (2010), N4, 9 pp.
\end{thebibliography}
\end{document}
